% Native semantic reconstruction of source physical page 28 (printed page 27).
% Provenance: IEEE Std 1516-2025, Clause 3.1, definitions subscribed attributes of an object class--time advancing state.
\HLAStartSemanticPage{28}{27}
\setlength{\HLASemanticDefinitionGap}{12.10pt}
\fontsize{10.02pt}{11.524pt}\selectfont
\begingroup
\hbadness=2500\relax
\addfontfeatures{FakeStretch=1.035}
\enlargethispage{12pt}
\vspace*{2.95pt}

\begin{hladefinition}[0pt]{def-subscribed-attributes-of-an-object-class}{subscribed attributes of an object class}
The class attributes that have been arguments to \HLAService{Subscribe Object Class Attributes} service invocations by a given joined federate for a given object class that have not subsequently been unsubscribed, either individually or by unsubscribing the whole object class.
\HLADefinitionNote[11.58pt]{See 5.1.2 and 5.8 of IEEE Std 1516.1\HLAStandardMark-2025.}
\end{hladefinition}

\begin{hladefinition}{def-subscribed-attributes-of-an-object-class-with-regions}{subscribed attributes of an object class with regions}
The class attributes that have been arguments to \HLAService{Subscribe Object Class Attributes With Regions} service invocations by a given joined federate for a given object class and a given region, assuming the joined federate did not subsequently invoke the \HLAService{Unsubscribe Object Class Attributes With Regions} service for that object class and region.
\HLADefinitionNote[11.58pt]{See 9.8 of IEEE Std 1516.1\HLAStandardMark-2025.}
\end{hladefinition}

\begin{hladefinition}{def-subscribed-interaction-class}{subscribed interaction class}
An interaction class that, from the perspective of a given joined federate, was an argument to a \HLAService{Subscribe Interaction Class} or \HLAService{Subscribe Interaction Class With Regions} service invocation that was not subsequently followed by an \HLAService{Unsubscribe Interaction Class} or \HLAService{Unsubscribe Interaction Class With Regions} service invocation for that interaction class.
\HLADefinitionNote[11.58pt]{See 5.1.3 and 5.10 of IEEE Std 1516.1\HLAStandardMark-2025.}
\end{hladefinition}

\begin{hladefinition}{def-subscribed-interaction-class-with-regions}{subscribed interaction class with regions}
An interaction class and a region that, from the perspective of a given joined federate, were arguments to a \HLAService{Subscribe Interaction Class With Regions} service invocation that was not subsequently followed by an \HLAService{Unsubscribe Interaction Class With Regions} service invocation for that interaction class and that region.
\HLADefinitionNote[11.58pt]{See 9.10 of IEEE Std 1516.1\HLAStandardMark-2025.}
\end{hladefinition}

\begin{hladefinition}{def-subscription-region-set}{subscription region set}
A set of regions used for subscription of a class attribute or used for subscription of an interaction class. \HLASeeAlso{\HLATermReference{used for subscription of a class attribute}\HLASemicolon{} \HLATermReference{used for subscription of an interaction class}.}
\end{hladefinition}

\begin{hladefinition}{def-superclass}{superclass}
A class that generalizes a set of properties that may be inherited by more refined (i.e., detailed) versions of the class. In High Level Architecture (HLA) applications, a class may have at most one immediate superclass (i.e., can inherit only from a single class at the next higher level of the class hierarchy).
\end{hladefinition}

\begin{hladefinition}{def-synchronization-point}{synchronization point}
A logical point in the sequence of a federation execution that all joined federates forming a synchronization set for that point attempt to reach and by which, if they are successful, they synchronize their respective processing.
\end{hladefinition}

\begin{hladefinition}{def-target-object-instance}{target object instance}
An object instance for which a directed interaction is intended.
\end{hladefinition}

\begin{hladefinition}[15.16pt]{def-time-advancing-state}{time advancing state}
The interval between a joined federate's request to advance its time position and the corresponding grant by the runtime infrastructure (RTI). A joined federate may advance its time position only by requesting a time advancement from the RTI via one of the following services:
\begin{hladefinitionservicelist}
\item \HLAService{Time Advance Request}
\item \HLAService{Time Advance Request Available}
\item \HLAService{Next Message Request}
\item \HLAService{Next Message Request Available}
\item \HLAService{Flush Queue Request}
\end{hladefinitionservicelist}
\par\vspace*{5.02pt}
The joined federate's time position will not actually be advanced until the RTI responds with a \HLAService{Time Advance Grant} \HLAPrintersDagger{} or \HLAService{Flush Queue Grant} \HLAPrintersDagger{} service invocation at that joined federate.
\end{hladefinition}
\endgroup
\clearpage
